\section{Institutional Interpretation}
\label{sec:institutional}

The fixed-capacity assumption is a stylized representation of policy prioritization rather than a literal claim that governments face a one-unit accounting identity. Smart-specialization policy provides a direct institutional analogue: regions are expected to identify a small number of priority areas and concentrate public resources on related transformative activities \citep{ForayEichlerKeller2021}. The model abstracts from the administrative details of that prioritization and isolates the opportunity cost created when greater emphasis on one domain necessarily displaces attention or resources from another.

Cross-regional complementarity also has a direct policy analogue. The European Union's 2026 Interregional Innovation Investments Strand 1 explicitly supports interregional investments built on shared or complementary smart-specialization priorities with the objective of strengthening EU value chains \citep{EISMEA2026}. More broadly, the OECD's framework for place-based industrial policy emphasizes alignment between national sectoral priorities and local capabilities and the importance of coordinated action across levels of government \citep{OECD2025}. These policy frameworks motivate the possibility that one region's priority can raise the return to a different but complementary activity elsewhere.

The model's local-return motive is likewise consistent with observed subnational policy behavior. \citet{TianWangZhang2027} find that Chinese local governments offer deeper land-price discounts to industries aligned with existing local specialization; they also report weaker support for industries linked to dominant local firms through supply chains and for specialization patterns extending across jurisdictional boundaries. This evidence is consistent with jurisdiction-bound policy incentives, but the present paper does not estimate $\Delta$, $A$, $\alpha$, or the general capture share $\lambda$ from those data.

The policy implication is correspondingly limited. The model does not establish that centralization is optimal, because it does not include administrative costs, local information advantages, or endogenous intergovernmental institutions. It shows instead that coordination of priority choices can have value when complementary benefits cross jurisdictional boundaries and the jurisdiction undertaking the complementary activity captures only part of the value it creates. In that setting, coordination can improve the composition of a fixed policy portfolio even without increasing the total amount of public support.
